\clearpage
\begin{proofbox}
\textbf{方法一：初中轴对称与直角三角形}
\begin{center}
\begin{tikzpicture}[x=1cm,y=1cm]
\draw[p055 frame] (0,0) rectangle (13.8,6.9);
\node[anchor=west,font=\small\bfseries] at (.4,6.4) {方法一：看直角三角形的边怎样翻折};
\begin{scope}[shift={(1.6,3.35)},scale=.95]
  \draw[p055 mirror] (-.65,0)--(4.7,0) node[right] {$y=0$};
  \coordinate (U) at (0,1); \coordinate (V) at (3,2);
  \coordinate (Up) at (0,-1); \coordinate (Vp) at (3,-2);
  \draw[p055 original] (U)--(V);
  \draw[p055 image] (Up)--(Vp);
  \draw[p055 helper] (U)--(3,1)--(V) (Up)--(3,-1)--(Vp);
  \draw[gray] (2.82,1)--(2.82,1.18)--(3,1.18);
  \draw[gray] (2.82,-1)--(2.82,-1.18)--(3,-1.18);
  \node[left] at (U) {$U$}; \node[above right] at (V) {$V$};
  \node[left] at (Up) {$U'$}; \node[below right] at (Vp) {$V'$};
  \node[below] at (1.5,1) {$\Delta x=3$};
  \node[right] at (3,1.5) {$\Delta y=1$};
  \node[above] at (1.5,-1) {$\Delta x'=3$};
  \node[right] at (3,-1.5) {$\Delta y'=-1$};
\end{scope}
\begin{scope}[shift={(8.65,1.45)},scale=1.05]
  \draw[p055 mirror] (0,0)--(3.5,3.5) node[above] {$y=x$};
  \coordinate (U) at (0,0); \coordinate (V) at (3,1); \coordinate (Vp) at (1,3);
  \draw[p055 original] (U)--(V); \draw[p055 image] (U)--(Vp);
  \draw[p055 helper] (U)--(3,0)--(V) (U)--(0,3)--(Vp);
  \draw[gray] (2.82,0)--(2.82,.18)--(3,.18);
  \draw[gray] (0,2.82)--(.18,2.82)--(.18,3);
  \node[below left] at (U) {$U=U'$};
  \node[right] at (V) {$V(3,1)$};
  \node[above left] at (Vp) {$V'(1,3)$};
  \node[below] at (1.5,0) {$3$}; \node[right] at (3,.5) {$1$};
  \node[left] at (0,1.5) {$3$}; \node[below] at (.5,3) {$1$};
\end{scope}
\node at (3.6,.55) {水平边同向，竖直边反向};
\node at (10.3,.55) {水平、竖直两条直角边互换};
\end{tikzpicture}

\end{center}
在原线上取不同两点 $U,V$，坐标差为 $\Delta x,\Delta y$。图中的水平、竖直直角边显示坐标差的绝对值；斜率还须保留方向的正负号。

\medskip
\textbf{水平、竖直镜面：} 关于 $y=b$，点变换为
\[
(x,y)\mapsto(x,2b-y),\qquad
(\Delta x',\Delta y')=(\Delta x,-\Delta y).
\]
关于 $x=a$，点变换为
\[
(x,y)\mapsto(2a-x,y),\qquad
(\Delta x',\Delta y')=(-\Delta x,\Delta y).
\]
因此当 $\Delta x\ne0$ 时，两种镜面都给出 $k'=-k$。

\medskip
\textbf{镜面 $y=x+c$：} 先向下平移 $c$，镜面成为 $y=x$。$(s,t)$ 与 $(t,s)$ 的中点在 $y=x$ 上，连线与镜面垂直，所以它们互为对称点（点在轴上时保持不动）。平移回去得到
\[
(x,y)\mapsto(y-c,x+c),\qquad
(\Delta x',\Delta y')=(\Delta y,\Delta x).
\]
图右侧的两条直角边因此互换，$k'=\Delta x/\Delta y=1/k$。

\medskip
\textbf{镜面 $y=-x+c$：} 对称点为 $(c-y,c-x)$，因为两点中点满足 $x+y=c$；不在镜面上时，连接两点的方向为 $(1,1)$，垂直于镜面；在镜面上时点保持不动。故增量变成 $(-\Delta y,-\Delta x)$，商仍为 $1/k$。

两种倒数规则都要求 $\Delta x\Delta y\ne0$；若其中一个增量为零，直接判断水平、竖直互换。
\end{proofbox}

\clearpage
\begin{proofbox}
\textbf{方法二：只跟踪方向增量}
\begin{center}
\begin{tikzpicture}[x=1cm,y=1cm]
\draw[p055 frame] (0,0) rectangle (13.8,10.6);
\node[anchor=west,font=\small\bfseries] at (.4,10.1) {方法二：只跟踪一根方向箭头（图取 $p=2,\ q=1$）};
\foreach \xx/\yy/\kind in {3.5/7.8/1,10.2/7.8/2,3.5/2.9/3,10.2/2.9/4}{
\begin{scope}[shift={(\xx,\yy)}]
  \draw[p055 axis] (-2.6,0)--(2.6,0);
  \draw[p055 axis] (0,-1.65)--(0,1.65);
  \draw[p055 original,-Stealth] (0,0)--(2,1);
  \node[CStatement,above] at (2,1) {$(p,q)$};
  \ifnum\kind=1
    \draw[p055 mirror] (0,-1.65)--(0,1.65);
    \draw[p055 image,-Stealth] (0,0)--(-2,1);
    \node[orange!85!black,above] at (-2,1) {$(-p,q)$};
    \node at (0,-1.95) {$x=a:\quad k'=-q/p$};
  \fi
  \ifnum\kind=2
    \draw[p055 mirror] (-2.6,0)--(2.6,0);
    \draw[p055 image,-Stealth] (0,0)--(2,-1);
    \node[orange!85!black,below] at (2,-1) {$(p,-q)$};
    \node at (0,-1.95) {$y=b:\quad k'=-q/p$};
  \fi
  \ifnum\kind=3
    \draw[p055 mirror] (-1.5,-1.5)--(1.65,1.65);
    \draw[p055 image,-Stealth] (0,0)--(1,2);
    \node[orange!85!black,above left] at (1,2) {$(q,p)$};
    \node at (0,-2.3) {$y=x+c:\quad k'=p/q$};
  \fi
  \ifnum\kind=4
    \draw[p055 mirror] (-1.65,1.65)--(1.5,-1.5);
    \draw[p055 image,-Stealth] (0,0)--(-1,-2);
    \node[orange!85!black,left] at (-1.1,-1.3) {$(-q,-p)$};
    \node at (0,-2.3) {$y=-x+c:\quad k'=p/q$};
  \fi
\end{scope}}
\end{tikzpicture}

\end{center}
\textbf{图怎样读：} 蓝色箭头是原增量 $(p,q)$，橙色箭头是像增量；细灰线只作横、纵方向参照。各小图把镜面上的一点移到局部原点，未把 $a,b,c$ 假设为零。平移不会改变方向增量。

\medskip
由轴对称对横向、纵向分量的作用（方法一已由点的反射验证），直接读得
\[
\begin{array}{c|c|c}
\text{镜面}&\text{像增量}&k'\\ \hline
x=a&(-p,q)&-q/p\\
y=b&(p,-q)&-q/p\\
y=x+c&(q,p)&p/q\\
y=-x+c&(-q,-p)&p/q
\end{array}
\]
原斜率为 $q/p$，所以前两类变号，后两类取倒数。这里不必计算任何点的实际位置。若像横增量为零，先报告像直线竖直；若原横增量为零，也应直接用增量判断，不写 $k$。
\end{proofbox}

\clearpage
\begin{proofbox}
\textbf{方法三：高中三角函数处理任意镜面}
\begin{center}
% alpha=30 degrees, theta=-15 degrees, theta'=75 degrees.
\begin{tikzpicture}[x=1cm,y=1cm]
\draw[p055 frame] (0,0) rectangle (13.8,7.6);
\node[anchor=west,font=\small\bfseries] at (.4,7.1) {方法三：镜面两侧等角，所以 $\theta'=2\alpha-\theta$};
\begin{scope}[shift={(3.1,2.55)}]
  \draw[p055 axis] (-1,0)--(6.5,0) node[right] {$x$ 轴正向};
  \draw[p055 mirror] (210:1.2)--(30:6.1) node[right] {$l$};
  \draw[p055 original,-Stealth] (0,0)--(-15:5.6) node[right] {$r$};
  \draw[p055 image,-Stealth] (0,0)--(75:3.6) node[above] {$r'$};
  \draw[CStatement,-Stealth] (0:1.2) arc (0:-15:1.2);
  \node[CStatement] at (-15:2.05) [below] {$\theta$};
  \draw[teal!70!black] (0:2.6) arc (0:30:2.6);
  \node[teal!70!black] at (15:3) {$\alpha$};
  \draw[orange!85!black] (0:3.6) arc (0:75:3.6);
  \node[orange!85!black] at (56:4.05) {$\theta'$};
  \draw[thick] (-15:.7) arc (-15:30:.7);
  \draw[thick] (30:.85) arc (30:75:.85);
  \node at (7.5,2.3) {$\alpha-\theta=\theta'-\alpha$};
  \node[align=center] at (7.5,1.2) {黑色小圆弧\\表示两个夹角相等};
  \fill (0,0) circle (1.8pt);
  \node[anchor=east] at (-.4,.6) {$P\in l$};
\end{scope}
\node at (6.9,.45) {图取 $\alpha=30^\circ,\ \theta=-15^\circ,\ \theta'=75^\circ$；一般方向允许相差 $180^\circ$。};
\end{tikzpicture}

\end{center}
设镜面 $y=mx+c$ 的方向角为 $\alpha$，原方向角为 $\theta$，像方向角为 $\theta'$，其中 $\tan\alpha=m,\ \tan\theta=k$。图中两个黑色小圆弧相等，因此
\[
\theta'=2\alpha-\theta\pmod{180^\circ}.
\]
可取 $\cos\alpha=1/\sqrt{1+m^2}$、$\sin\alpha=m/\sqrt{1+m^2}$，同样取 $\cos\theta=1/\sqrt{1+k^2}$、$\sin\theta=k/\sqrt{1+k^2}$。于是
\[
\sin2\alpha=\frac{2m}{1+m^2},\qquad
\cos2\alpha=\frac{1-m^2}{1+m^2}.
\]
若像方向不竖直，差角公式给出
\[
\begin{aligned}
k'&=\frac{\sin2\alpha\cos\theta-\cos2\alpha\sin\theta}
{\cos2\alpha\cos\theta+\sin2\alpha\sin\theta}\\[3pt]
&=\boxed{\frac{2m+(m^2-1)k}{1-m^2+2mk}}.
\end{aligned}
\]
分母为零正好对应 $\cos(2\alpha-\theta)=0$，所以像直线竖直。$m=0$ 给出 $-k$；$m=\pm1$ 且 $k\ne0$ 给出 $1/k$。镜面或原线竖直时，不套有限斜率通式，改用方向判断。
\end{proofbox}

\clearpage
\begin{proofbox}
\textbf{方法四：垂足与中点的坐标法}
\begin{center}
\begin{tikzpicture}[x=1cm,y=1cm]
\draw[p055 frame] (0,0) rectangle (13.8,7.6);
\node[anchor=west,font=\small\bfseries] at (.4,7.1) {方法四：垂足 $H_i$ 是每一对 $Q_i,Q_i'$ 的中点};
\begin{scope}[shift={(2.1,2.9)},scale=.92]
  \coordinate (Q1) at (0,3); \coordinate (Q2) at (2,3.5);
  \coordinate (H1) at (1.2,.6); \coordinate (H2) at (3,1.5);
  \coordinate (R1) at (2.4,-1.8); \coordinate (R2) at (4,-.5);
  \draw[p055 mirror] (-.8,-.4)--(6,3) node[right] {$l$};
  \draw[p055 original,-Stealth] (Q1)--(Q2);
  \draw[p055 image,-Stealth] (R1)--(R2);
  \draw[p055 helper] (Q1)--(R1) (Q2)--(R2);
  \draw[gray] ($(H1)+(.2,.1)$)--($(H1)+(.1,.3)$)--($(H1)+(-.1,.2)$);
  \draw[gray] ($(H2)+(.2,.1)$)--($(H2)+(.1,.3)$)--($(H2)+(-.1,.2)$);
  \foreach \u/\v in {Q1/H1,H1/R1}{
    \draw[gray] ($ (\u)!.5!(\v) +(-.13,-.065)$)--($ (\u)!.5!(\v) +(.13,.065)$);
  }
  % Distinct tick codes: only each reflected pair has equal halves.
  \foreach \u/\v in {Q2/H2,H2/R2}{
    \foreach \t in {.46,.54}{
      \draw[gray] ($ (\u)!\t!(\v) +(-.13,-.065)$)--($ (\u)!\t!(\v) +(.13,.065)$);
    }
  }
  \foreach \p/\label/\pos in {Q1/Q_1/above left,Q2/Q_2/above,R1/Q_1'/left,R2/Q_2'/right}{
    \fill (\p) circle (1.8pt) node[\pos] {$\label$};
  }
  \fill (H1) circle (1.8pt); \node at (.4,1) {$H_1$};
  \fill (H2) circle (1.8pt); \node at (3.3,2.2) {$H_2$};
  \node[CStatement,anchor=west] at (7.2,3.5) {原方向 $(1,k)$};
  \node[orange!85!black,below right] at (3.4,-1.2) {$(\Delta x',\Delta y')$};
  \node[align=left,anchor=west] at (7.2,1.8) {$Q_iH_i\perp l$\\[5pt]$Q_iH_i=H_iQ_i'$\\[5pt]$Q_i'=2H_i-Q_i$};
\end{scope}
\node at (6.9,.45) {图取 $l:y=\frac12x$；虚线是垂线，短刻线表示两半等长。};
\end{tikzpicture}

\end{center}
\textbf{先求一个点：} 对 $Q(x,y)$ 与镜面 $y=mx+c$，垂线方向可取 $(-m,1)$，它与镜面方向 $(1,m)$ 的数量积为 $0$。设垂足 $H=(x-mt,y+t)$，代入镜面得
\[
t=\frac{mx+c-y}{1+m^2},\qquad Q'=2H-Q=(x-2mt,y+2t).
\]
记 $d=y-mx-c$，则
\[
Q'=\left(x+\frac{2md}{1+m^2},\ y-\frac{2d}{1+m^2}\right).
\]
\textbf{再对两个点相减：} 在原线上取 $Q_1,Q_2$，使 $x_2-x_1=1,\ y_2-y_1=k$。此时 $d_2-d_1=k-m$，所以
\[
\begin{aligned}
\Delta x'&=1+\frac{2m(k-m)}{1+m^2}=\frac{1-m^2+2mk}{1+m^2},\\
\Delta y'&=k-\frac{2(k-m)}{1+m^2}=\frac{2m+(m^2-1)k}{1+m^2}.
\end{aligned}
\]
横增量不为零时，两式相除即得通式；为零时像直线竖直。反射保持两点距离，故两个像增量不可能同时为零。$c$ 在相减时消去，只说明方向与截距无关，不能据此忽略位置。
\end{proofbox}
